\documentclass[11pt,a4paper]{article}
\usepackage[utf8]{inputenc}
\usepackage[vietnamese]{babel}
\usepackage{amsmath,amssymb}
\usepackage{geometry}
\geometry{a4paper, total={180mm,260mm}, left=15mm, right=15mm, top=18mm, bottom=18mm}

\usepackage{tikz}
\usetikzlibrary{shapes.geometric, arrows.meta, positioning, calc}
\tikzset{
  arr/.style={-{Stealth[scale=0.8]}},
  every node/.append style={align=center}
}
\usepackage{tcolorbox}
\tcbuselibrary{skins,breakable}
\usepackage{fontawesome}
\usepackage{fancyhdr}
\usepackage{enumitem}

\pagestyle{fancy}
\fancyhf{}
\renewcommand{\headrulewidth}{0.5pt}
\renewcommand{\footrulewidth}{0.5pt}
\fancyhead[L]{\small\bfseries\color{violet!85!black}\faBook\ SỔ TAY KIẾN THỨC TOÁN 2}
\fancyhead[R]{\small\bfseries\color{violet!85!black}Lớp Toán Thầy Lĩnh -- loptoanthaylinh.vn}
\fancyfoot[C]{\small\color{gray}-- Sổ tay kiến thức: Bài 19. Phép cộng có nhớ số có 2 chữ số với số có 1 chữ số --}

\setlength{\parindent}{0pt}
\setlength{\parskip}{0pt}

\begin{document}

% --- BANNER BÀI HỌC ---
\begin{tcolorbox}[
  enhanced,
  colback=violet!8!white,
  colframe=violet!70!black,
  arc=3mm, boxrule=1.2pt,
  left=3mm, right=3mm, top=2mm, bottom=2mm,
  halign=center
]
  {\LARGE\bfseries\color{violet!90!black}\faBookmark\ BÀI 19: PHÉP CỘNG (CÓ NHỚ) SỐ CÓ HAI CHỮ SỐ VỚI SỐ CÓ MỘT CHỮ SỐ}\\[1mm]
  {\itshape\footnotesize\color{black!75} Kĩ thuật đặt tính thẳng cột và nhớ 1 sang hàng chục dạng 35 + 7 trên 1 trang A4}
\end{tcolorbox}

\vspace{-2mm}

\noindent
\begin{minipage}[t]{0.49\textwidth}\vspace{0pt}
\begin{tcolorbox}[
  enhanced,
  colback=violet!4!white,
  colframe=violet!75!black,
  arc=2mm, boxrule=0.8pt,
  left=2mm, right=2mm, top=1.8mm, bottom=1.8mm,
  title={\footnotesize\bfseries\color{white}\faKey\ 1. Quy tắc đặt tính}
]
{\footnotesize
\begin{itemize}\setlength{\itemsep}{1pt}\setlength{\leftskip}{-2.5mm}
  \item \textbf{Đặt tính thẳng cột:}\\
  - Viết số có hai chữ số ở trên.\\
  - Viết số có một chữ số ở dưới sao cho thẳng cột với chữ số \textbf{hàng đơn vị}.\\
  - Viết dấu cộng ($+$) và kẻ vạch ngang.
  \item \textbf{Thực hiện tính từ phải sang trái:}
  \begin{itemize}\setlength{\itemsep}{0.5pt}\setlength{\leftskip}{-2.5mm}
    \item Cộng chữ số hàng đơn vị với nhau.
    \item Nếu tổng $\ge 10$, viết chữ số đơn vị và \textbf{nhớ 1 sang hàng chục}.
    \item Chữ số hàng chục \textbf{thêm 1 vừa nhớ} rồi viết vào kết quả.
  \end{itemize}
\end{itemize}
}
\end{tcolorbox}
\end{minipage}
\hfill
\begin{minipage}[t]{0.49\textwidth}\vspace{0pt}
\begin{tcolorbox}[
  enhanced,
  colback=teal!4!white,
  colframe=teal!75!black,
  arc=2mm, boxrule=0.8pt,
  left=2mm, right=2mm, top=1.8mm, bottom=1.8mm,
  title={\footnotesize\bfseries\color{white}\faSliders\ 2. Phân tích chi tiết ví dụ $35 + 7$}
]
{\footnotesize
\begin{center}
\begin{tabular}{r@{\,}ll}
  $3$ & $5$ & \quad \textbf{Thao tác tính:} \\
$+$ & & \quad $\bullet$ $5 + 7 = 12$, \\
    & $7$ & \quad \ \ viết $\mathbf{2}$, \textbf{nhớ 1}. \\ \cline{1-2}
  $\mathbf{4}$ & $\mathbf{2}$ & \quad $\bullet$ $3$ \textbf{thêm 1 bằng 4}, viết $\mathbf{4}$.
\end{tabular}
\end{center}
\vspace{1mm}
$\implies$ Kết quả: $35 + 7 = \mathbf{42}$.\\
\textbf{Nhận xét:} Hàng chục từ $3$ tăng lên thành $4$ do có nhớ $1$.
}
\end{tcolorbox}
\end{minipage}

\vspace{-1.5mm}
\begin{tcolorbox}[
  enhanced,
  colback=orange!4!white,
  colframe=orange!80!black,
  arc=2mm, boxrule=0.8pt,
  left=2mm, right=2mm, top=1.5mm, bottom=1.5mm,
  title={\footnotesize\bfseries\color{white}\faSitemap\ 3. Sơ đồ tiến trình thao tác tính có nhớ}
]
\centering
\begin{tikzpicture}[x=1.3cm, y=0.8cm]
  \node[rectangle, rounded corners=2pt, draw=violet!80!black, fill=violet!15, font=\scriptsize\bfseries] (b1) at (0,0) {HÀNG ĐƠN VỊ\\$5 + 7 = 12$\\Viết $\mathbf{2}$};
  \draw[->, very thick, red!85!black] (b1) -- node[above, font=\tiny\bfseries] {Nhớ 1 chục} (2.8,0);
  \node[rectangle, rounded corners=2pt, draw=teal!80!black, fill=teal!15, font=\scriptsize\bfseries] (b2) at (4.5,0) {HÀNG CHỤC\\$3$ thêm $1 = \mathbf{4}$\\Viết $\mathbf{4}$};
  \draw[->, very thick, green!70!black] (b2) -- node[above, font=\tiny\bfseries] {Kết quả} (6.8,0);
  \node[rectangle, rounded corners=2pt, draw=green!80!black, fill=green!20, font=\normalsize\bfseries] at (7.6,0) {$\mathbf{42}$};
\end{tikzpicture}
\end{tcolorbox}

\vspace{-1.5mm}
\noindent
\begin{minipage}[t]{0.49\textwidth}\vspace{0pt}
\begin{tcolorbox}[
  enhanced,
  colback=green!4!white,
  colframe=green!60!black,
  arc=2mm, boxrule=0.8pt,
  left=2mm, right=2mm, top=1.8mm, bottom=1.8mm,
  title={\footnotesize\bfseries\color{white}\faCheckCircle\ 4. Ví dụ thực hành chuẩn mực}
]
{\footnotesize
\textbf{Đặt tính rồi tính:}\\
\begin{tabular}{r@{\,}lr@{\,}lr@{\,}l}
  $4$ & $8$ & $6$ & $4$ & $7$ & $9$ \\
$+$ & & $+$ & & $+$ & \\
    & $5$ & & $8$ & & $3$ \\ \hline
  $\mathbf{5}$ & $\mathbf{3}$ & $\mathbf{7}$ & $\mathbf{2}$ & $\mathbf{8}$ & $\mathbf{2}$
\end{tabular}
\quad
\begin{tabular}{l}
  $\bullet$ $8+5=13$, viết $3$ nhớ $1$; $4+1=5$.\\
  $\bullet$ $4+8=12$, viết $2$ nhớ $1$; $6+1=7$.\\
  $\bullet$ $9+3=12$, viết $2$ nhớ $1$; $7+1=8$.
\end{tabular}
}
\end{tcolorbox}
\end{minipage}
\hfill
\begin{minipage}[t]{0.49\textwidth}\vspace{0pt}
\begin{tcolorbox}[
  enhanced,
  colback=red!4!white,
  colframe=red!75!black,
  arc=2mm, boxrule=0.8pt,
  left=2mm, right=2mm, top=1.8mm, bottom=1.8mm,
  title={\footnotesize\bfseries\color{white}\faExclamationTriangle\ 5. Mẹo tránh bẫy trắc nghiệm}
]
{\footnotesize
\begin{itemize}\setlength{\itemsep}{1pt}\setlength{\leftskip}{-2.5mm}
  \item \textbf{Lỗi quên cộng nhớ:}\\
  $35 + 7 = 32$ (quên không thêm 1 vào 3) là lỗi phổ biến nhất. Hãy chấm 1 chấm nhỏ trên đầu chữ số hàng chục để nhắc nhớ.
  \item \textbf{Lỗi đặt số 1 chữ số lệch sang cột chục:}\\
  Viết $7$ dưới $3$ $\implies$ sai hoàn toàn.
\end{itemize}
}
\end{tcolorbox}
\end{minipage}


\end{document}
